\documentclass[10pt,a4paper]{article}
\usepackage[T1]{fontenc}
\usepackage[utf8]{inputenc}
\usepackage{lmodern}
\usepackage[a4paper,top=1.6cm,bottom=1.6cm,left=2.1cm,right=2.1cm]{geometry}
\usepackage{microtype}
\usepackage{enumitem}
\usepackage[hidelinks]{hyperref}
\setlength{\parindent}{0pt}
\setlength{\parskip}{0.55em}
\pagestyle{empty}
\newcommand{\letterhead}{%
\begin{flushright}
\textbf{Dr.\ Kundai Farai Sachikonye}\\
AIMe Registry for Artificial Intelligence\\
Email: \href{mailto:kundai.sachikonye@bitspark.com}{kundai.sachikonye@bitspark.com}\\
Phone: +49 162 6386344
\end{flushright}
\vspace{0.4em}
5 October 2026\par
\vspace{0.4em}
The Editor\\
Editorial Office\par}
\newcommand{\declarations}{%
\textbf{Declarations.}
\begin{itemize}[leftmargin=1.2em,itemsep=0.15em,topsep=0.2em]
\item The manuscript is original, has not been published previously, and is not under consideration for publication elsewhere.
\item I am the sole author, have approved the submitted version, and take full responsibility for its content.
\item I declare no competing financial or non-financial interests.
\item This work received no specific grant from any funding agency in the public, commercial, or not-for-profit sectors.
\end{itemize}}
\newcommand{\signoff}{%
Thank you for considering this manuscript. I look forward to the reviewers' assessment and am happy to supply any further information the editor requires.

Yours sincerely,\par
\vspace{0.8em}
Dr.\ Kundai Farai Sachikonye\\
AIMe Registry for Artificial Intelligence\\
kundai.sachikonye@bitspark.com \quad|\quad +49 162 6386344}

\begin{document}
\letterhead

\textbf{Re: Submission of manuscript, ``Motor Commands and Proprioceptive Return as Phases of a Single Charge Circulation: Resolving the Deafferentation Paradox''}

Dear Editor,

I am pleased to submit the above manuscript for consideration as an original research article in your journal.

Standard motor control theory treats the nervous system as an open-loop controller corrected by sensory feedback. That picture predicts that losing proprioception should degrade the precision of movement. The clinical record shows something different: in deafferented patients, of whom Ian Waterman is the most extensively studied, and in the Taub--Berman--Mott primate series, coordinated movement is abolished outright even though motor output and intent are intact. This dissociation has remained unexplained for three decades.

The manuscript resolves it by showing that motor output and proprioceptive return are not two signals coupled by feedback but two phases of one closed charge circulation. Starting from three accepted foundations (the Hodgkin--Huxley membrane circuit, sliding-filament contraction, and Kirchhoff's laws applied to the complete neural--muscular--proprioceptive network), it derives the musculoskeletal system as a closed circuit with no external ground. Charge conservation together with variance minimisation then requires every motor command to complete its circulation through the proprioceptive return. The consequences follow directly:
\begin{itemize}[leftmargin=1.2em,itemsep=0.25em]
\item deafferentation prevents movement completion, because the circuit cannot close, rather than merely degrading movement;
\item the stretch reflex is an autocatalytic oscillator, not a negative-feedback servo;
\item postural control is continuous oscillation, not dynamic equilibrium;
\item the locomotor central pattern generator is a self-sustaining closed-loop resonance in the spinal cord;
\item motor learning is a geometric modification of circuit conductance, not the storage of commands.
\end{itemize}

The manuscript gives formal derivations and quantitative predictions for stretch-reflex latency (30--50~ms), the force--length and force--velocity relations, the postural sway power spectrum, Henneman's size principle, and the effect of deafferentation on movement generation. These are checked in six numerical experiments, including a Hodgkin--Huxley simulation and a closed-loop postural pendulum compared against its deafferented counterpart. The manuscript then places the framework alongside the equilibrium-point hypothesis, optimal feedback control, and internal-model theories. It shows that each of these is a partial description within the closed-circuit picture, and that only the closed-circuit picture accounts for the deafferentation paradox. It also sets out empirical tests that distinguish it from these alternatives, together with implications for neural prostheses and robotics. I believe the work will interest readers in motor control, neurophysiology, rehabilitation, and biomechanics.

\declarations

\signoff
\end{document}
